\section{Discussion: What This Means for an Investment Decision}
\label{sec:discussion}

This section is written for a reader who needs to decide whether to act
on this project's output, not for a reader who needs to audit its
mathematics.

\subsection{What the model actually tells you}

The headline finding of this project is simple: a computer program,
given only historical prices and a handful of publicly available
industry statistics -- and given \emph{no} instruction about which of
those statistics matter -- consistently figures out on its own that
crude-oil price swings are tied most closely to \emph{physical,
industrial} activity (how much the oil-and-gas-support sector of the
economy is growing or shrinking) and to the \emph{other} major
benchmark crude price (Brent), and only loosely, if at all, to broad
financial-market sentiment (the stock market, the dollar, bond yields).
That is a genuinely useful piece of information for anyone trying to
understand \emph{why} oil prices move, independent of whether the model
can also tell you \emph{which direction} they will move next.

\subsection{What the model does not (yet) tell you}

Separately, and just as important: when we asked the same model to make
a concrete, one-month-ahead call on whether oil would be in a rising,
falling, or flat period, it was right about 7 times out of 10 -- but
simply assuming "next month looks like this month" was right closer to
9 times out of 10. In plain terms: on this test, the model's
market-timing signal was not yet good enough to act on by itself. This
is a common and important finding in quantitative investing more
generally -- many signals that look statistically interesting
completely fail to beat the simplest possible comparison (just assuming
things stay the same), and the discipline of checking against that
simple comparison, rather than skipping straight to "the model found
something", is exactly what separates a credible research process from
an overconfident one.

\subsection{What we would recommend to a decision-maker}

\begin{itemize}
  \item \textbf{Use the model's structural finding as a monitoring
    tool, not a trading signal.} Knowing that industrial activity and
    the Brent benchmark are the variables to watch is actionable today,
    at no risk: it tells a risk manager or a policy analyst \emph{where}
    to focus limited attention, even without trusting the model's
    specific price-direction calls.
  \item \textbf{Do not allocate real capital based on the current
    forecast alone.} A trading or hedging decision sized directly off
    this model's monthly regime call would, on the evidence gathered
    here, likely underperform the much simpler rule of just assuming
    "no change." Before any capital is put behind this specific
    forecast, it would need to clear the same bar in a longer,
    independent test period, and ideally be combined with other signals
    rather than used alone.
  \item \textbf{The portfolio-tilting demonstration (Section
    \ref{sec:results}) illustrates a mechanism, not a recommendation.}
    We showed that a confident model view \emph{can} be translated into
    a concrete portfolio tilt (shifting an equal-weighted position
    toward oil when the model is bullish). That machinery is sound and
    reusable, but plugging in today's specific, not-yet-validated
    forecast would be putting a mechanism ahead of the evidence that
    mechanism should be fed by.
  \item \textbf{Cost-effectiveness.} Every input used by this project
    (Yahoo Finance prices, FRED and EIA public statistics) is free and
    automatically refreshable, and the entire pipeline runs in minutes
    on an ordinary laptop. The honest 70\%-vs-88\% comparison above cost
    nothing extra to produce and should be considered a required,
    essentially free quality check before any such model is trusted with
    a real decision -- there is no excuse, cost-wise, for skipping it.
\end{itemize}

In short: this project demonstrates a transparent, inexpensive, and
genuinely interpretable way to monitor what is driving the oil market,
and a clean mechanism for turning a model view into a portfolio
decision -- but it also demonstrates, honestly, that the specific
market-timing signal produced by this particular model, on this
particular sample, is not yet strong enough to deploy on its own.
